\section*{CHƯƠNG 3. KẾ HOẠCH QUẢN LÝ PHẠM VI DỰ ÁN}
\addcontentsline{toc}{section}{CHƯƠNG 3. KẾ HOẠCH QUẢN LÝ PHẠM VI DỰ ÁN}

Chương này xác lập văn bản \textit{Scope Statement} (Tuyên bố phạm vi dự án) nhằm định rõ quy cách sản phẩm chuyển giao và ranh giới công việc, đồng thời thiết lập \textit{Work Breakdown Structure -- WBS} (Cấu trúc phân rã công việc) kèm từ điển mô tả chi tiết và quy trình kiểm soát thay đổi phạm vi.

\subsection*{3.1. Tuyên bố phạm vi dự án và đặc tả sản phẩm}
\addcontentsline{toc}{subsection}{3.1. Tuyên bố phạm vi dự án và đặc tả sản phẩm}

Văn bản Tuyên bố phạm vi dự án là cam kết chính thức giữa Ban Quản lý dự án, Phòng Giáo dục và Đào tạo cùng năm trường trung học cơ sở thụ hưởng về các sản phẩm bàn giao.

Sản phẩm chuyển giao cuối cùng là 2.000 cuốn ấn phẩm ``Bộ tài liệu tư vấn phân luồng và định hướng nghề nghiệp cho học sinh cuối cấp trung học cơ sở'' kèm theo 01 bộ tệp tin chế bản điện tử gốc. Ấn phẩm được thiết kế theo khổ tiêu chuẩn 17 $\times$ 24 cm, gồm \textbf{80 trang ruột} (chia thành 5 tay sách 16 trang, in trên giấy Fort trắng định lượng 100 gsm chống lóa) và \textbf{4 trang bìa ngoài} (Bìa 1, 2, 3, 4 in trên giấy Couche định lượng 250 gsm cán màng mờ mặt ngoài), in offset bốn màu toàn bộ và gia công đóng gáy keo nhiệt.

Nội dung 80 trang ruột được cấu trúc thành phần mở đầu và bốn phần chuyên môn (tổng cộng 20 bài chuẩn), kế thừa quy trình định hướng nghề nghiệp bốn bước của Lê Thị Quỳnh Nga, Nguyễn Thị Hương và Chu Cẩm Thơ (2026) cùng nghiên cứu chuẩn hóa công cụ mật mã Holland của Phạm Mạnh Hà (2020):

Trang đầu sách (từ trang 1 đến trang 4, gồm 4 trang): Trang phụ bìa, lời giới thiệu, hướng dẫn cách khai thác tài liệu và mục lục tổng thể.

Phần 1 -- Khám phá tiềm năng và nhận diện bản thân (từ trang 5 đến trang 22, gồm 18 trang, 5 bài chuẩn): Giới thiệu bài trắc nghiệm sở thích nghề nghiệp theo mật mã Holland rút gọn gồm 36 câu hỏi dành cho lứa tuổi 14 đến 15, hướng dẫn tự chấm điểm, vẽ biểu đồ nhận diện nhóm sở thích nổi trội và đối chiếu với năng lực học tập các môn văn hóa.

Phần 2 -- Bản đồ các con đường học tập sau tốt nghiệp trung học cơ sở (từ trang 23 đến trang 44, gồm 22 trang, 6 bài chuẩn): Hệ thống hóa bằng sơ đồ đồ họa thông tin và bảng so sánh bốn hướng đi sau lớp 9 (trường trung học phổ thông công lập, trường trung học phổ thông ngoài công lập, trung tâm Giáo dục nghề nghiệp -- Giáo dục thường xuyên và các trường trung cấp, cao đẳng đào tạo hệ chín cộng), làm rõ điều kiện xét tuyển, thời gian đào tạo và lộ trình học liên thông.

Phần 3 -- Thế giới nghề nghiệp và chính sách hỗ trợ học nghề (từ trang 45 đến trang 62, gồm 18 trang, 5 bài chuẩn): Cung cấp thông tin dự báo nhu cầu nhân lực tại địa phương và giải thích cụ thể quy định miễn 100\% học phí học nghề đối với học sinh tốt nghiệp trung học cơ sở theo Nghị định số 81/2021/NĐ-CP và Nghị định số 97/2023/NĐ-CP.

Phần 4 -- Cẩm nang ra quyết định và đồng hành cùng gia đình (từ trang 63 đến trang 80, gồm 18 trang, 4 bài chuẩn): Hướng dẫn quy trình năm bước lựa chọn nguyện vọng, cung cấp bộ phiếu bài tập đối thoại giữa cha mẹ và học sinh, kèm danh bạ thông tin tuyển sinh của 25 cơ sở giáo dục nghề nghiệp uy tín trên địa bàn.

\subsection*{3.2. Cấu trúc phân rã công việc}
\addcontentsline{toc}{subsection}{3.2. Cấu trúc phân rã công việc}

\subsubsection*{3.2.1. Phương pháp tiếp cận và nguyên lý xây dựng}

Dự án áp dụng phương pháp \textit{Top-down} (Tiếp cận từ trên xuống) kết hợp kỹ thuật \textit{Mind-mapping} (Sơ đồ tư duy) để phân rã tổng thể dự án thành 5 giai đoạn ở cấp độ 1 và 20 gói công việc cụ thể ở cấp độ 2.

Cấu trúc phân rã công việc tuân thủ bảy nguyên lý chuẩn mực: bao quát đúng 100\% phạm vi công việc đã cam kết; mỗi gói công việc chỉ xuất hiện tại một vị trí duy nhất; mỗi gói công việc ở cấp thấp nhất chỉ gắn với một sản phẩm đầu ra và do duy nhất một cá nhân chịu trách nhiệm chính; phù hợp với trình tự tổ chức thực hiện thực tế; có sự tham gia xây dựng của các tiểu ban chuyên môn và được định nghĩa rõ trong bảng từ điển công việc.

\subsubsection*{3.2.2. Sơ đồ cây phân rã công việc}

\begin{figure}[htbp]
\centering
\caption{Sơ đồ cây phân rã công việc (WBS) của dự án}
\label{fig:so_do_wbs}
\resizebox{0.95\textwidth}{!}{%
\begin{tikzpicture}[
    node distance = 0.4cm and 0.3cm,
    root_box/.style = {rectangle, rounded corners=3pt, draw=red!80!black, fill=red!10, line width=1.1pt, text width=13cm, align=center, inner sep=5pt, font=\footnotesize\bfseries},
    level1_box/.style = {rectangle, rounded corners=3pt, draw=blue!80!black, fill=blue!10, line width=0.9pt, text width=2.4cm, align=center, inner sep=4pt, font=\scriptsize\bfseries},
    level2_box/.style = {rectangle, rounded corners=2pt, draw=gray!70, fill=gray!6, line width=0.6pt, text width=2.4cm, align=left, inner sep=3pt, font=\tiny},
    arr/.style = {-{Stealth[scale=0.8]}, line width=0.7pt, draw=gray!70!black}
]

% Root
\node (root) [root_box] {DỰ ÁN XÂY DỰNG BỘ TÀI LIỆU TƯ VẤN PHÂN LUỒNG TRUNG HỌC CƠ SỞ};

% Level 1
\node (g1) [level1_box, below left=0.8cm and 3.5cm of root] {1.0 KHẢO SÁT \&\\ĐỀ CƯƠNG};
\node (g2) [level1_box, right=0.35cm of g1] {2.0 BIÊN SOẠN\\BẢN THẢO};
\node (g3) [level1_box, right=0.35cm of g2] {3.0 THẨM ĐỊNH\\\& HIỆU ĐÍNH};
\node (g4) [level1_box, right=0.35cm of g3] {4.0 CHẾ BẢN\\\& IN ẤN};
\node (g5) [level1_box, right=0.35cm of g4] {5.0 KIỂM ĐỊNH\\\& BÀN GIAO};

\draw [arr] (root.south) -- ++(0,-0.3) -| (g1.north);
\draw [arr] (root.south) -- ++(0,-0.3) -| (g2.north);
\draw [arr] (root.south) -- ++(0,-0.3) -| (g3.north);
\draw [arr] (root.south) -- ++(0,-0.3) -| (g4.north);
\draw [arr] (root.south) -- ++(0,-0.3) -| (g5.north);

% Level 2 - Branch 1
\node (g1_1) [level2_box, below=0.3cm of g1] {1.1 Khảo sát nhu cầu (350 phiếu)};
\node (g1_2) [level2_box, below=0.15cm of g1_1] {1.2 Dữ liệu và phân tích số liệu};
\node (g1_3) [level2_box, below=0.15cm of g1_2] {1.3 Xây dựng, duyệt đề cương};
\draw [arr] (g1.south) -- (g1_1.north);

% Level 2 - Branch 2
\node (g2_1) [level2_box, below=0.3cm of g2] {2.1 Viết Phần 1 (18 trang)};
\node (g2_2) [level2_box, below=0.15cm of g2_1] {2.2 Viết Phần 2 (22 trang)};
\node (g2_3) [level2_box, below=0.15cm of g2_2] {2.3 Viết Phần 3 (18 trang)};
\node (g2_4) [level2_box, below=0.15cm of g2_3] {2.4 Viết Phần 4 (18 trang)};
\node (g2_5) [level2_box, below=0.15cm of g2_4] {2.5 Biên tập bản thảo lần 1};
\draw [arr] (g2.south) -- (g2_1.north);

% Level 2 - Branch 3
\node (g3_1) [level2_box, below=0.3cm of g3] {3.1 Phản biện 3 chuyên gia};
\node (g3_2) [level2_box, below=0.15cm of g3_1] {3.2 Họp Hội đồng 2 phiên};
\node (g3_3) [level2_box, below=0.15cm of g3_2] {3.3 Lấy ý kiến đọc thử};
\node (g3_4) [level2_box, below=0.15cm of g3_3] {3.4 Hiệu đính bản sạch};
\draw [arr] (g3.south) -- (g3_1.north);

% Level 2 - Branch 4
\node (g4_1) [level2_box, below=0.3cm of g4] {4.1 Lựa chọn nhà thầu in};
\node (g4_2) [level2_box, below=0.15cm of g4_1] {4.2 Vẽ hình, dàn trang};
\node (g4_3) [level2_box, below=0.15cm of g4_2] {4.3 In thử duyệt bản bông};
\node (g4_4) [level2_box, below=0.15cm of g4_3] {4.4 In offset 2.000 cuốn};
\draw [arr] (g4.south) -- (g4_1.north);

% Level 2 - Branch 5
\node (g5_1) [level2_box, below=0.3cm of g5] {5.1 Kiểm định mẫu (384 c)};
\node (g5_2) [level2_box, below=0.15cm of g5_1] {5.2 Họp Hội đồng nghiệm thu};
\node (g5_3) [level2_box, below=0.15cm of g5_2] {5.3 Đóng thùng, bàn giao};
\node (g5_4) [level2_box, below=0.15cm of g5_3] {5.4 Thanh lý \& lưu hồ sơ};
\draw [arr] (g5.south) -- (g5_1.north);

\end{tikzpicture}%
}
\end{figure}

\subsubsection*{3.2.3. Bảng từ điển phân rã công việc chi tiết}

Tài liệu \textit{WBS Dictionary} (Từ điển phân rã công việc) tại Bảng \ref{tab:tu_dien_wbs} xác định rõ phạm vi tác nghiệp, kết quả bàn giao và cá nhân chịu trách nhiệm chính cho toàn bộ 20 gói công việc, làm cơ sở trực tiếp để lập dự toán chi phí tại Chương 5.

\begin{table}[htbp]
\centering
\caption{Từ điển phân rã công việc chi tiết của dự án}
\label{tab:tu_dien_wbs}
\begin{tabularx}{\textwidth}{cCYYL{3.4cm}}
\toprule
\textbf{Mã} & \textbf{Tên gói công việc} & \textbf{Nội dung triển khai} & \textbf{Sản phẩm đầu ra} & \textbf{Chịu trách nhiệm} \\
\midrule
\textbf{1.0} & \multicolumn{3}{l}{\textbf{Khảo sát nhu cầu và lập khung cấu trúc}} & \textbf{PGS. TS. Nguyễn Minh Hà} \\
1.1 & Khảo sát thực tế & Thu thập 350 bộ phiếu khảo sát (300 HS, 50 GV/PH) tại 5 trường & 350 bộ phiếu hợp lệ & Chuyên viên điều tra \\
1.2 & Dữ liệu tuyển sinh & Xử lý số liệu, tổng hợp chỉ tiêu 25 trường nghề và chính sách học phí & Báo cáo kết quả dữ liệu & Thư ký dự án \\
1.3 & Duyệt đề cương & Thiết kế ma trận 4 phần và hội thảo chuyên môn (12 đại biểu) & Bản đề cương được duyệt & Chủ biên chuyên môn \\
\midrule
\textbf{2.0} & \multicolumn{3}{l}{\textbf{Biên soạn bản thảo bộ tài liệu (80 trang ruột)}} & \textbf{PGS. TS. Nguyễn Minh Hà} \\
2.1 & Viết Phần 1 & Viết 5 bài và chuẩn hóa bảng hỏi Holland 36 câu (trang 5--22) & Bản thảo P1 (18 trang) & Trưởng nhóm TG1 \\
2.2 & Viết Phần 2 & Viết 6 bài và lập ma trận so sánh 4 hướng đi sau THCS (trang 23--44) & Bản thảo P2 (22 trang) & Trưởng nhóm TG2 \\
2.3 & Viết Phần 3 & Viết 5 bài xu hướng nghề và chính sách miễn 100\% học phí (trang 45--62) & Bản thảo P3 (18 trang) & Trưởng nhóm TG3 \\
2.4 & Viết Phần 4 & Viết 4 bài, phiếu đối thoại gia đình và danh bạ 25 trường (trang 63--80) & Bản thảo P4 (18 trang) & Trưởng nhóm TG4 \\
2.5 & Biên tập bản thảo 1 & Soạn 4 trang đầu sách, ráp nối 20 bài và biên tập 80 trang ruột & Bản thảo lần một (80 tr) & Chủ biên chuyên môn \\
\midrule
\textbf{3.0} & \multicolumn{3}{l}{\textbf{Thẩm định và hiệu đính bản thảo}} & \textbf{Tiến sĩ Vũ Thị Mai} \\
3.1 & Phản biện độc lập & 3 chuyên gia độc lập đọc phản biện toàn văn 80 trang bản thảo & 3 bản nhận xét phản biện & Chủ tịch HĐ thẩm định \\
3.2 & Họp Hội đồng & Tổ chức 2 phiên họp Hội đồng thẩm định (5 thành viên) & Biên bản họp Hội đồng & Chủ tịch HĐ thẩm định \\
3.3 & Lấy ý kiến đọc thử & Thảo luận nhóm tập trung đọc thử với 20 học sinh và 5 giáo viên & Báo cáo ý kiến đọc thử & Thư ký HĐ thẩm định \\
3.4 & Hiệu đính bản sạch & Bốn nhóm tác giả chỉnh sửa theo kết luận, rà soát chính tả & Bản thảo sạch hoàn chỉnh & Chủ biên chuyên môn \\
\midrule
\textbf{4.0} & \multicolumn{3}{l}{\textbf{Thiết kế chế bản và in ấn tài liệu}} & \textbf{ThS. Phạm Quang Thắng} \\
4.1 & Lựa chọn nhà thầu & Chào hàng cạnh tranh rút gọn và ký hợp đồng trọn gói & Báo cáo thầu và Hợp đồng & Phụ trách Mua sắm \\
4.2 & Thiết kế dàn trang & Thiết kế 4 trang bìa ngoài, vẽ 25 hình minh họa và dàn trang điện tử & Tệp ma-két hoàn chỉnh & Nhóm thiết kế chế bản \\
4.3 & In thử duyệt mẫu & In thử 2 bộ bản bông màu, rà soát lỗi và ký duyệt niêm phong & Bản bông có chữ ký duyệt & Chủ biên chuyên môn \\
4.4 & In offset hàng loạt & Bình bản xuất kẽm, in 2.000 cuốn, cán màng mờ và đóng keo nhiệt & 2.000 cuốn tài liệu & Giám đốc nhà thầu in \\
\midrule
\textbf{5.0} & \multicolumn{3}{l}{\textbf{Kiểm định, nghiệm thu và bàn giao}} & \textbf{Thạc sĩ Lê Hoàng Nam} \\
5.1 & Kiểm định mẫu & Rút mẫu phân tầng 384 cuốn kiểm tra quy cách kỹ thuật tại kho & Phiếu kiểm định 384 cuốn & Đại diện đơn vị thụ hưởng \\
5.2 & Họp nghiệm thu & Hội đồng nghiệm thu cấp cơ sở đánh giá sản phẩm và ký biên bản & Biên bản nghiệm thu & Giám đốc dự án \\
5.3 & Đóng thùng bàn giao & Đóng 40 thùng giấy (50 c/thùng), vận chuyển bàn giao cho 5 trường & Biên bản bàn giao 5 trường & Đại diện đơn vị thụ hưởng \\
5.4 & Thanh lý \& lưu trữ & Quyết toán hợp đồng in ấn, đóng gói 03 bộ hồ sơ lưu trữ gốc & Biên bản thanh lý \& hồ sơ & Phụ trách Mua sắm \\
\bottomrule
\end{tabularx}
\end{table}

\subsection*{3.3. Quy trình xác nhận và kiểm soát sự thay đổi phạm vi}
\addcontentsline{toc}{subsection}{3.3. Quy trình xác nhận và kiểm soát sự thay đổi phạm vi}

Xác nhận phạm vi được thực hiện tại điểm kết thúc của từng giai đoạn thông qua việc đối chiếu sản phẩm trung gian thực tế với tiêu chí đầu ra trong Bảng \ref{tab:tu_dien_wbs}. Nhằm ngăn ngừa hiện tượng \textit{Scope Creep} (Trượt phạm vi dự án), Ban Quản lý dự án áp dụng hai chốt kiểm soát định lượng: khống chế dung lượng phần ruột ở mức cố định 80 trang (5 tay sách) để không làm biến động đơn giá chế bản và giấy in; đồng thời từ chối mọi đề xuất bổ sung các hoạt động ngoài phạm vi xuất bản học liệu đã được phê duyệt tại Điều lệ dự án.

